\section[Persistent VWF forests]{Persistent exact VWF forest elimination and explicit ranges}
\label{sec:op3-persistent-vwf-forest}

\paragraph{Status and source comparison.}
CPW Lemma 6.1 and Lemma 6.6, PDF pp.\ 30--34
\citep{chen2021diffusion}, already state fast degree-one elimination using
maintained vertex functions. The source was visually checked, including
Algorithm 6's elimination-before-compression schedule. We do not claim
degree-one elimination as a new algorithmic idea. This section gives a
separately implemented derivative-integral representation, direct proofs
of its exact constants and all allocated work, and explicit bounds for one
uncompressed forest pass. These are \statusproved{} drafts awaiting review.
The scalar coefficient sign on source pp.\ 40--42 is corrected below.
No fast coarse oracle or recursive complexity theorem is imported here.

\begin{lemma}[Correct Lift coefficients and signed endpoints; \statusproved]
\label{lem:op3-correct-vwf-lift}
Let $F$ be a finite-piece VWF on $[L,\infty)$, $L\leq0$, and $c>0$.
For a parent endpoint $P\leq0$ define
$\widehat F(x)=\min_{y\geq L}\{c(x-y)^2/2+F(y)\}$ on $[P,\infty)$.
Then $\widehat F$ is a VWF with the same final derivative as $F$.
If the minimizing $y$ lies in a piece
$F(y)=ry^2/2+ay+b$, the correct new piece is
\begin{equation}
 \widehat F(x)=\frac{cr}{2(c+r)}x^2+
 \frac{ca}{c+r}x+b-\frac{a^2}{2(c+r)}.
 \label{eq:op3-correct-vwf-lift-piece}
\end{equation}
The boundary transition is $s_0=L+F'(L)/c$. For $x\leq s_0$ the
minimizer is $L$ and the function is $F(L)+c(x-L)^2/2$.
Every old split $s$ maps to $s+F'(s)/c$ with derivative value $F'(s)$.
For all $x\geq P$ the unique minimizer is
\begin{equation}
 y(x)=x-\widehat F'(x)/c.
 \label{eq:op3-vwf-lift-recovery}
\end{equation}
In particular, if $P>s_0$, the exact new endpoint value is
$F(y(P))+\widehat F'(P)^2/(2c)$; otherwise it is $F(L)+c(P-L)^2/2$.
\end{lemma}
\begin{proof}
The function minimized in $y$ is strongly convex and grows quadratically
above a global affine lower bound for $F$, so it has a unique minimizer.
Its interior stationarity equation gives $y=(cx-a)/(c+r)$.
Substitution gives \eqref{eq:op3-correct-vwf-lift-piece}, with a minus
sign in its constant. At the boundary the one-sided stationarity condition
is $c(L-x)+F'(L)\geq0$, equivalent to $x\leq s_0$.
The split map is strictly increasing because $F'$ is nondecreasing.
Interior curvature changes from $r$ to $cr/(c+r)$, a nondecreasing
function of $r$ bounded by $c$. Thus the decreasing sequence of old
curvatures remains decreasing, including the initial boundary curvature
$c$ and the final zero curvature. Values and derivatives match at every
transition, as direct substitution shows. Restricting to $[P,\infty)$
discards only an initial prefix of pieces. Moreover
$\widehat F(0)\leq F(0)\leq0$ by the feasible choice $y=0$.
Envelope differentiation gives \eqref{eq:op3-vwf-lift-recovery}, including
the boundary branch, and proves the endpoint formulas.
\end{proof}

\paragraph{Source formula diagnostic.}
The printed positive constant correction on source pp.\ 40--42 disagrees
with \eqref{eq:op3-correct-vwf-lift-piece}. The valid VWF $F(y)=-y$ on
$y\geq0$, with $c=1$, has $\widehat F(0)=-1/2$ at $y=1$; a positive
correction gives $+1/2$. This is a specific coefficient correction, not
a refutation of the whole source theorem. Recovery here explicitly uses
the saved lifted child function. Source Algorithm 7, p.\ 33, instead
prints the parent descriptor $D_v$ after its Add operation in the recovery
line, whereas Lemma 6.6 defines OptimalX after Lift. Our reconstruction
contract is \eqref{eq:op3-vwf-lift-recovery}, avoiding that descriptor
ambiguity and keeping each child's lifted snapshot independently available.

\begin{lemma}[Affine updates of derivative integrals; \statusproved]
\label{lem:op3-affine-derivative-integrals}
Represent a piecewise linear derivative by ordered points $(x,g)$.
A subtree stores its first and last point and the trapezoid integral
$I=\int g\,dx$ between them. Under an order-preserving affine point map
\[
 x'=ax+bg+t_x,\qquad g'=cx+dg+t_g,
\]
the aggregate changes in constant exact-real word work to
\begin{equation}
 I'=(ad-bc)I+bc\,\Delta(xg)+\frac{ac}{2}\Delta(x^2)
       +\frac{bd}{2}\Delta(g^2)+t_g\bigl(a\Delta x+b\Delta g\bigr),
 \label{eq:op3-affine-derivative-area}
\end{equation}
where each $\Delta$ is last endpoint minus first endpoint.
An immutable AVL tree with these aggregates supports a derivative query
or prefix integral in $O(\log(2+k))$ work without allocating tree nodes,
where $k$ is the number of stored points.
\end{lemma}
\begin{proof}
Expand $\int(cx+dg+t_g)(a\,dx+b\,dg)$ and use
$\int x\,dg=\Delta(xg)-\int g\,dx$ and the endpoint identities for
$\int x\,dx$ and $\int g\,dg$. This gives
\eqref{eq:op3-affine-derivative-area}; translation $t_x$ disappears from
the differential. The formula also holds on every individual segment.

A node combines the two child aggregates and the connecting trapezoids.
A lazy affine update transforms the node point and aggregate, composes its
pending tag, and shares its unchanged children. It does not reconstruct
an aggregate using children still in their old coordinates. A read-only
query carries the composed map down one search path. Fully covered child
integrals are included in constant work, with one partial connecting
segment at the query endpoint. The AVL height recurrence bounds the path
by $O(\log(2+k))$; no push or path copy is needed for reads.
\end{proof}

\paragraph{Implemented scalar operations.}
The new module \texttt{persistent\_vwf\_forest.py} preserves the old
validated affine-tree and removal modules. Its augmented node stores
the domain point $(L,F'(L))$, all retained derivative knots, and the
value $F(L)$ separately. The final derivative is constant. Function values
use the prefix integral, with the constant derivative on the final ray.

Lift applies $(x,g)\mapsto(x+g/c,g)$ to one root. Before its transformed
domain point, it uses derivative $c(x-L)$. At parent endpoint $P$, it
computes the exact constant by \cref{lem:op3-correct-vwf-lift}, deletes
the transformed points at or below $P$ when necessary, and inserts the
new endpoint. Repeated minimum deletion is used; its total cost is charged
below, rather than assumed to be one logarithmic operation.

For addition on a common domain $[P,\infty)$, the smaller derivative has
the form $g(P)+C(x-P)-\sum_s q_s(x-s)_+$, with $q_s\geq0$.
One global vertical shear adds the affine term to the larger tree.
A stream through the smaller tree then inserts each needed split and
applies its negative suffix hinge. The temporary final curvature is tracked
until it returns to zero. It stays nonnegative throughout: it is the
remaining curvature of the smaller function, added to the larger function's
nonnegative curvature. Zero drops are skipped. Constants at $P$ are added.
The operations preserve earlier snapshots by copying only touched paths.
Canonical export streams positive curvature drops, omits zero jumps, and
uses the exact $F(0),F'(0)$ anchors. It costs linear work in the curve size
and returns the atom representation of \cref{sec:op3-global-vwf-compression}.

\begin{theorem}[Implemented persistent supplied-forest elimination; \statusproved]
\label{thm:op3-persistent-vwf-forest}
Suppose a supplied weighted graph with $n$ numbered vertices and $m$ edges
is the union of vertex-disjoint rooted trees and edges joining their $r$
retained roots. Every tree has exactly one retained root. Give vertex $i$
a VWF $f_i$ on $[L_i,\infty)$, $L_i\leq0$, with $k_i$ splits, and set
$W=\sum_i(k_i+2)$. The implementation validates the topology and produces
exact retained-root functions, the copied root-to-root edge list, and an
immutable reconstruction representation in
\begin{equation}
 O\bigl(m+W\log^2(2+W)\bigr)\ \text{word work},\qquad
 O\bigl(m_R+W\log^2(2+W)\bigr)\ \text{newly allocated words},
 \label{eq:op3-persistent-vwf-forest-work}
\end{equation}
where $m_R$ is the root-edge count. These bounds include all copies and
retained versions, not only peak live storage. The retained functions have
total size $O(W)$ and can be exported in $O(W)$ additional work and words.

For any feasible supplied root vector, one downward pass costs
$O(n\log(2+W))$ work and total newly allocated words, including its
$O(n)$ output vector, and returns the exact
conditional minimum over all eliminated coordinates. Its original energy
equals the reduced root energy. The original root gradients equal the
reduced gradients. A single-root tree with nonnegative total final slope
can therefore be minimized exactly with one additional derivative-zero
query. No unknown-support graph discovery is included in these supplied
access bounds.
\end{theorem}
\begin{proof}
Orient the supplied forest from its given roots, reading each incidence
and storing every parent edge weight. Reject extra incidences at eliminated
vertices. Copy each root edge once. A second incidence pass constructs the
postorder merge inputs. All orientation arrays, descriptors, unsuccessful
topology checks and copied edge records are charged in $O(n+m)$ work and
$O(n+m_R)$ words. Vertices are array-indexed; reconstruction reads stored
weights and makes no edge-key searches.

For a vertex $v$, define bottom-up
\[
 F_v=f_v+\sum_{w:\,p(w)=v}\widehat F_w,\qquad
 \widehat F_w(x)=\min_{y\geq L_w}
       \{c_{wv}(x-y)^2/2+F_w(y)\}.
\]
Each lifted child is restricted to the parent's domain. The scalar lemma
proves these identities exactly. Induction shows that $F_v$ is the
subtree conditional minimum for a fixed value at $v$. Given parent value
$x$, the saved $\widehat F_w'$ yields the unique minimizing child value
by \eqref{eq:op3-vwf-lift-recovery}. A downward pass therefore realizes
all conditional minima simultaneously. Only the retained root edges remain
outside those subtrees, proving energy equality. Differentiating the lifted
contributions gives the root-gradient equality. Finite single-root minima
and zero total final slope are covered by
\cref{prop:op3-vwf-canonical-input-gap}; the root derivative is continuous
nondecreasing with a nonnegative final value, so a zero crossing or its
domain endpoint is sufficient.

It remains to account for all representation work. Assign each original
vertex immutable mass $k_i+2$, and each subtree the sum of its input masses.
At a vertex choose the largest-mass input among its own function and all
lifted children as the shared base. Every other input has mass at most
that base's mass. Its point count is at most its mass: original knots,
domain points and at most one newly created endpoint per eliminated vertex
are enough to account for every retained record. Each smaller input's
stream can thus be charged to its original mass tokens. Whenever a token
is charged, its containing merged mass at least doubles. There are at most
$O(\log(2+W))$ charges per token, for $O(W\log(2+W))$ streamed points.
Every streamed hinge makes $O(\log(2+W))$ path visits and node copies.
Local input construction is linear in its points; selecting merge inputs
and retaining all response pointers is linear in the graph and input size.

Prefix deletion does not escape this ledger. Consider only the logical
current curves on the postorder frontier; saved inactive snapshots are
never reintroduced into that frontier. Initially there are
$\sum_i(k_i+1)$ points. A merge consumes the smaller curve and creates
at most its non-domain point count in the base, so it cannot increase the
frontier's total point count. Each Lift introduces at most one endpoint.
Consequently the total number of points removed by prefix deletion is at
most $\sum_i(k_i+1)+n-r=O(W)$. Each removal costs logarithmic path work
and allocation. Each remaining Lift uses one lazy root update, logarithmic
endpoint queries and at most one insertion. These facts establish
\eqref{eq:op3-persistent-vwf-forest-work}, including temporary paths,
rotations, affine tags, integral aggregates, output records and all old
snapshots. Each read-only derivative query in reconstruction is logarithmic;
its temporary carry records are charged as well as its output vector.
Exporting all root curves is
linear because their disjoint mass sum is $W$.
\end{proof}

\begin{proposition}[One-pass terminal, curvature and split bounds; \statusproved]
\label{prop:op3-vwf-forest-ranges}
In \cref{thm:op3-persistent-vwf-forest}, let $t_i$ be the final derivative
of $f_i$, $h_i$ its largest curvature, $S_+=\sum_i\max(t_i,0)$,
$U=\max(0,\text{all input splits})$, and let every forest edge have
weight at least $c_{\min}>0$. For a subtree rooted at $v$, put
$T_v=\sum_{i\text{ in that subtree}}t_i$. Then
\begin{enumerate}
\item $F_v$ and its lifted version have final derivative $T_v$, including
when it is negative.
\item Every retained split in the exact pass is at most
$U+nS_+/c_{\min}$ and at least the lower endpoint of its current domain.
Thus all split magnitudes are bounded by
$\max\{\max_i(-L_i),U+nS_+/c_{\min}\}$.
\item The largest curvature of $F_v$ is at most
$h_v+\sum_{w:p(w)=v}c_{vw}$, and that of $\widehat F_v$ is at most
$c_{vp(v)}$. In particular the sum of largest curvatures of all retained
root functions is at most $\sum_i h_i+\sum_{e\text{ in forest}}c_e$.
\item If the full original objective has a finite minimum $\Phi^*$,
then every stored $F_v(0)$ and $\widehat F_v(0)$ belongs to $[\Phi^*,0]$.
\end{enumerate}
No intermediate compression is used in these claims, and no new graph
edge is produced by the pass.
\end{proposition}
\begin{proof}
Lift preserves the final derivative and addition sums it, proving the first
claim by induction. Every event or endpoint $s$ lifted from a subtree maps
to $s+F_v'(s)/c$, which is at most
$\max(0,s)+\max(T_v,0)/c\leq\max(0,s)+S_+/c_{\min}$.
Original domain endpoints are nonpositive. New parent endpoints are also
nonpositive, and prefix restriction discards events rather than moving them
farther. An event traverses at most $n-1$ forest edges. This proves the
split bound, including signed input splits and the boundary transition.

The Lift curvature formula is at most its edge weight. Adding the children
to $f_v$ gives the curvature bound, and summing over roots is bounded by
the larger sum over all vertices. These operations alter only vertex
functions and remove forest edges; the retained graph edges are copied
unchanged.

Choosing zero throughout a subtree gives nonpositive conditional energy,
so both kinds of stored functions are nonpositive at zero. For a lower
bound, embed a conditional minimizer with root value zero into the full
graph by setting all outside coordinates to zero. The boundary forest edge
has zero contribution in the $F_v(0)$ case; it is already part of the
minimum in the $\widehat F_v(0)$ case. There are no other incidences out
of the eliminated subtree. The remaining vertex constants are nonpositive,
and all other edge contributions are zero. The full feasible energy is
therefore no greater than the stored value at zero and is at least
$\Phi^*$. Conditional minimizers exist by repeated strongly convex scalar
Lift. This proves the last claim.
\end{proof}

\begin{corollary}[Terminal mass of a generic proximal instance; \statusproved]
\label{cor:op3-generic-proximal-terminal-mass}
For graph Laplacians in \cref{thm:op3-generic-relative-proximal}, let
$W_G$ be the sum of the supplied $\mathcal G$ edge weights and let
$S_+=\sum_i\max(T_i,0)$ for the original vertex final derivatives.
At any accelerated center, the normalized model's positive final-slope
mass is at most
\[
 S_++32\kappa\sqrt{\kappa W_G\Delta}.
\]
This bounds the input $S_+$ for a subsequent valid forest pass without
requiring coordinate bounds or nonnegative individual shifted final slopes.
\end{corollary}
\begin{proof}
For any graph Laplacian $L$ with edge-weight sum $W_L$,
$\|Lx\|_1\leq2\sum_e c_e|x_i-x_j|
\leq2\sqrt{W_L}\|x\|_L$.
Trace comparison gives $W_H\leq\kappa W_G$. Since
$\|x\|_G\leq\|x\|_H\leq8\kappa\sqrt\Delta$, the triangle
inequality yields
$\|(\mathcal G-\mathcal H)x\|_1
\leq32\kappa\sqrt{\kappa W_G\Delta}$.
The normalized model changes each final slope by
$((\mathcal G-\mathcal H)x)_i$. The positive-part mass can increase by
at most the $\ell_1$ norm of this shift. This argument does not assume
edgewise nonnegative weights for the difference $\mathcal H-\mathcal G$.
\end{proof}

\paragraph{Audit and limits.}
The explicit reference audit checks 3,000 scalar Lift cases, 16,997 complete
constant/recovery polynomial identities, and 192 tree reconstructions.
The persistent audit independently checks 1,000 Lift/addition/persistence
cases, 34,404 complete derivative/quadratic overlay intervals, 5,941 saved
curve versions, 9,081 general affine segment/prefix identities and all
stored AVL/integral aggregates. It reconstructs 192 atlas trees, certifies
15 retained-root field assignments, rejects three invalid elimination
topologies, and checks original KKT on nine path/star/binary instances
through 512 vertices. All coordinates in the nine larger solutions are
positive. Canonical export is compared with full polynomial coefficients.
Exact reference matrices are unnecessary for this primitive: independent
scalar pieces and original sparse KKT equations supply the validators.

The immutable affine-tree and removal backends retain their prior hashes.
No floating-point stability or rational bit-complexity bound is claimed.
The remaining bridge is an implemented compressed coarse oracle with a
compatible additive-error budget, followed by the actual recursive graph
and preconditioner schedule. The forest bound on $c_{\min}$ is an input
hypothesis; it does not prove a lower bound for future recursively built
graphs. A positive compression error still needs a justified energy-scale
budget. Unknown-support discovery and cumulative graph-local work remain
separate. General OP3 remains \statusopen{}.
